% GENERATED FILE. Edit canonical lessons, then run npm run generate:books.
% canonical-source-hash: fnv1a64:041f133923d084e0
% canonical-lessons: BN-C305-prodip, BN-C305-atoshbaji, BN-C305-alpona, BN-C305-jholmol-kora, BN-C305-phanus

\chapter{Lights and Fireworks}
\label{ch:bn-lights-and-fireworks}
\hlchaptermodality{305}

\begin{chapteropening}
\textbf{By the end of this chapter:} \emph{I can describe festival lights, floor paintings and fireworks.}

Complete the last lesson of chapter 305: I can describe festival lights, floor paintings and fireworks.
\end{chapteropening}

\section[\texorpdfstring{\bn{প্রদীপ} (prôdīp)}{প্রদীপ (prôdīp)}]{\bn{প্রদীপ} (prôdīp) — an oil lamp}

\label{lesson:BN-C305-prodip}



\begin{quote}
Before the new one: say the Bengali for a pandal, then the Bengali for to celebrate.
\end{quote}

\subsection*{You'll want to know: \bn{প্রদীপ}}
\textbf{\bn{প্রদীপ}} — \emph{prôdīp} — \textquotedblleft{}an oil lamp\textquotedblright{}.

\begin{quote}
\bn{সন্ধ্যায়} \bn{দিদিমা} \bn{বারান্দায়} \bn{একটা} \bn{প্রদীপ} \bn{জ্বালালেন।}
\end{quote}

\begin{quote}
\emph{sôndhyāy didimā bārāndāy ekṭā prôdīp jvālālen.}
\end{quote}

\begin{quote}
In the evening my grandmother lit an oil lamp on the veranda.
\end{quote}

\begin{cousinweb}
\textbf{\bn{প্রদীপ}} (\emph{prôdīp}) is Sanskrit \emph{pradīpa}, 'a light, a lamp' (Monier-Williams), from \emph{pra-√dīp}, 'to blaze forth'.
\end{cousinweb}

\begin{grammarlens}[title={What you've built}]
One more word for telling what happened.
\end{grammarlens}

\subsection*{Your turn}
\begin{itemize}
\raggedright
  \item \emph{Say it:} \emph{prôdīp}
  \item \emph{Say it:} \emph{prôdīp}, once more
  \item \emph{Say it:} \emph{sôndhyāy didimā bārāndāy ekṭā prôdīp jvālālen.} — In the evening my grandmother lit an oil lamp on the veranda.
  \item \emph{Recall:} say the Bengali for a pandal, then the Bengali for to celebrate, then say \emph{prôdīp} again
\end{itemize}

\begin{tcolorbox}[breakable,colback=teal!4,colframe=teal!35!black,title={Before you move on}]
What does \bn{প্রদীপ} mean? (\textquotedblleft{}An oil lamp\textquotedblright{}.) Say it once more.
\end{tcolorbox}
\section[\texorpdfstring{\bn{আতশবাজি} (ātôshbāji)}{আতশবাজি (ātôshbāji)}]{\bn{আতশবাজি} (ātôshbāji) — fireworks}

\label{lesson:BN-C305-atoshbaji}



\begin{quote}
Before the new one: say the Bengali for to celebrate, then the Bengali for an oil lamp.
\end{quote}

\subsection*{You'll want to know: \bn{আতশবাজি}}
\textbf{\bn{আতশবাজি}} — \emph{ātôshbāji} — \textquotedblleft{}fireworks\textquotedblright{}.

\begin{quote}
\bn{রাতে} \bn{আকাশে} \bn{আতশবাজি} \bn{জ্বলছিল।}
\end{quote}

\begin{quote}
\emph{rāte ākāshe ātôshbāji jvôlchhilô.}
\end{quote}

\begin{quote}
Fireworks were going off in the night sky.
\end{quote}

\begin{cousinweb}
\textbf{\bn{আতশবাজি}} (\emph{ātôshbāji}) is usually traced to Persian \emph{ātash-bāzī}: \emph{ātash}, 'fire', and \emph{bāzī}, 'play'.
\end{cousinweb}

\begin{grammarlens}[title={What you've built}]
One more word for telling what happened.
\end{grammarlens}

\subsection*{Your turn}
\begin{itemize}
\raggedright
  \item \emph{Say it:} \emph{ātôshbāji}
  \item \emph{Say it:} \emph{ātôshbāji}, once more
  \item \emph{Say it:} \emph{rāte ākāshe ātôshbāji jvôlchhilô.} — Fireworks were going off in the night sky.
  \item \emph{Recall:} say the Bengali for to celebrate, then the Bengali for an oil lamp, then say \emph{ātôshbāji} again
\end{itemize}

\begin{tcolorbox}[breakable,colback=teal!4,colframe=teal!35!black,title={Before you move on}]
What does \bn{আতশবাজি} mean? (\textquotedblleft{}Fireworks\textquotedblright{}.) Say it once more.
\end{tcolorbox}
\section[\texorpdfstring{\bn{আলপনা} (ālpônā)}{আলপনা (ālpônā)}]{\bn{আলপনা} (ālpônā) — a floor painting (alpona)}

\label{lesson:BN-C305-alpona}



\begin{quote}
Before the new one: say the Bengali for an oil lamp, then the Bengali for fireworks.
\end{quote}

\subsection*{You'll want to know: \bn{আলপনা}}
\textbf{\bn{আলপনা}} — \emph{ālpônā} — \textquotedblleft{}a floor painting (alpona)\textquotedblright{}.

\begin{quote}
\bn{পুজোর} \bn{আগে} \bn{মেয়েরা} \bn{দরজার} \bn{সামনে} \bn{আলপনা} \bn{আঁকল।}
\end{quote}

\begin{quote}
\emph{pujor āge meyerā dôrjār sāmne ālpônā ā̃klô.}
\end{quote}

\begin{quote}
Before the puja the girls painted an alpona in front of the door.
\end{quote}

\begin{cousinweb}
\textbf{\bn{আলপনা}} (\emph{ālpônā}) comes down from Sanskrit \emph{ālimpana}, 'whitening, painting' (CDIAL 1385, which lists Bengali \emph{ālpanā}): the white patterns drawn on a floor for a festival.
\end{cousinweb}

\begin{grammarlens}[title={What you've built}]
One more word for telling what happened.
\end{grammarlens}

\subsection*{Your turn}
\begin{itemize}
\raggedright
  \item \emph{Say it:} \emph{ālpônā}
  \item \emph{Say it:} \emph{ālpônā}, once more
  \item \emph{Say it:} \emph{pujor āge meyerā dôrjār sāmne ālpônā ā̃klô.} — Before the puja the girls painted an alpona in front of the door.
  \item \emph{Recall:} say the Bengali for an oil lamp, then the Bengali for fireworks, then say \emph{ālpônā} again
\end{itemize}

\begin{tcolorbox}[breakable,colback=teal!4,colframe=teal!35!black,title={Before you move on}]
What does \bn{আলপনা} mean? (\textquotedblleft{}A floor painting (alpona)\textquotedblright{}.) Say it once more.
\end{tcolorbox}
\section[\texorpdfstring{\bn{ঝলমল} \bn{করা} (jhôlmôl kôrā)}{ঝলমল করা (jhôlmôl kôrā)}]{\bn{ঝলমল} \bn{করা} (jhôlmôl kôrā) — to glitter; to sparkle}

\label{lesson:BN-C305-jholmol-kora}



\begin{quote}
Before the new one: say the Bengali for fireworks, then the Bengali for a floor painting.
\end{quote}

\subsection*{You'll want to know: \bn{ঝলমল} \bn{করা}}
\textbf{\bn{ঝলমল} \bn{করা}} — \emph{jhôlmôl kôrā} — \textquotedblleft{}to glitter; to sparkle\textquotedblright{}.

\begin{quote}
\bn{আলোয়} \bn{পুরো} \bn{শহরটা} \bn{ঝলমল} \bn{করছিল।}
\end{quote}

\begin{quote}
\emph{āloy puro shôhôrṭā jhôlmôl kôrchhilô.}
\end{quote}

\begin{quote}
The whole town was glittering with lights.
\end{quote}

Past (he/she …): \bn{সেই} \bn{শহর} \bn{আলোয়} \bn{ঝলমল} \bn{করল}, \emph{sei shôhôr āloy jhôlmôl kôrlô}.

\begin{cousinweb}
\textbf{\bn{ঝলমল} \bn{করা}} (\emph{jhôlmôl kôrā}) puts \bn{ঝলমল}, an expressive word whose doubled sound is the flicker of light, before \bn{করা}. No dictionary we use records an older form, so its history is not settled.
\end{cousinweb}

\begin{grammarlens}[title={What you've built}]
One more word for telling what happened.
\end{grammarlens}

\subsection*{Your turn}
\begin{itemize}
\raggedright
  \item \emph{Say it:} \emph{jhôlmôl kôrā}
  \item \emph{Say it:} \emph{jhôlmôl kôrā}, once more
  \item \emph{Say it:} \emph{āloy puro shôhôrṭā jhôlmôl kôrchhilô.} — The whole town was glittering with lights.
  \item \emph{Recall:} say the Bengali for fireworks, then the Bengali for a floor painting, then say \emph{jhôlmôl kôrā} again
\end{itemize}

\begin{tcolorbox}[breakable,colback=teal!4,colframe=teal!35!black,title={Before you move on}]
What does \bn{ঝলমল} \bn{করা} mean? (\textquotedblleft{}To glitter; to sparkle\textquotedblright{}.) Say it once more.
\end{tcolorbox}
\section[\texorpdfstring{\bn{ফানুস} (phānus)}{ফানুস (phānus)}]{\bn{ফানুস} (phānus) — a sky lantern}

\label{lesson:BN-C305-phanus}



\begin{quote}
Before the new one: say the Bengali for a floor painting, then the Bengali for to glitter.
\end{quote}

\subsection*{You'll want to know: \bn{ফানুস}}
\textbf{\bn{ফানুস}} — \emph{phānus} — \textquotedblleft{}a sky lantern\textquotedblright{}.

\begin{quote}
\bn{উৎসবের} \bn{রাতে} \bn{আকাশে} \bn{অনেক} \bn{ফানুস} \bn{উড়ছিল।}
\end{quote}

\begin{quote}
\emph{utsôber rāte ākāshe ônek phānus uṛchhilô.}
\end{quote}

\begin{quote}
On the night of the festival many sky lanterns were floating in the sky.
\end{quote}

\begin{cousinweb}
\textbf{\bn{ফানুস}} (\emph{phānus}) is usually traced to Persian \emph{fānūs}, 'a lantern'; in Bengal it is the paper lantern sent up into the sky.
\end{cousinweb}

\begin{grammarlens}[title={What you've built}]
One more word for telling what happened.
\end{grammarlens}

\subsection*{Your turn}
\begin{itemize}
\raggedright
  \item \emph{Say it:} \emph{phānus}
  \item \emph{Say it:} \emph{phānus}, once more
  \item \emph{Say it:} \emph{utsôber rāte ākāshe ônek phānus uṛchhilô.} — On the night of the festival many sky lanterns were floating in the sky.
  \item \emph{Recall:} say the Bengali for a floor painting, then the Bengali for to glitter, then say \emph{phānus} again
\end{itemize}

\begin{tcolorbox}[breakable,colback=teal!4,colframe=teal!35!black,title={Before you move on}]
What does \bn{ফানুস} mean? (\textquotedblleft{}A sky lantern\textquotedblright{}.) Say it once more.
\end{tcolorbox}
